% Generated draft: requires mathematical and source review.
\documentclass[11pt]{article}
\usepackage[margin=25mm]{geometry}
\usepackage{fontspec}
\setmainfont{Latin Modern Roman}
\usepackage{amsmath,amssymb,mathtools,unicode-math}
\setmathfont{Latin Modern Math}
\usepackage{enumitem,needspace,xurl,hyperref}
\setlength{\parindent}{0pt}
\setlength{\parskip}{5pt}
\setlength{\emergencystretch}{3em}
\newcommand{\N}{\mathbb{N}}
\newcommand{\Z}{\mathbb{Z}}
\newcommand{\Q}{\mathbb{Q}}
\newcommand{\R}{\mathbb{R}}
\newcommand{\C}{\mathbb{C}}
\newcommand{\F}{\mathbb{F}}
\newcommand{\A}{\mathbb{A}}
\newcommand{\PP}{\mathbb{P}}
\newcommand{\E}{\mathbb{E}}
\newcommand{\eps}{\varepsilon}
\newcommand{\dd}{\,\mathrm{d}}
\begin{document}
\section*{UnsolvedMath ID: 20002664}
\textbf{Source problem code:} catalogue entry
\subsection{Definitions and mathematical statement}
The registry fuses two consecutive Free Analysis questions: multivariable Cauchy/Poisson smoothing, and \(L^2\)-closure of cyclic gradients. The intervening ``1Q:'' is a page-number artifact, not one combined question.

For the first, the normalized scalar Poisson kernel is \(P_\varepsilon(z)=\varepsilon/[\pi(z^2+\varepsilon^2)]\), \(\varepsilon>0\). With a scalar Cauchy \(Z_\varepsilon\) independent of a Hermitian matrix \(A\), central noise \(C_\varepsilon=Z_\varepsilon I\) gives
\[
 \mathbb E_C[f(A+C_\varepsilon)\mid A]=(P_\varepsilon*f)(A)
\]
for bounded Borel \(f\). If \(A\) is random this is a conditional expectation over the added noise. The source's damaged kernel and unspecified ``matrix Cauchy'' law require review; central noise is one explicit realization, not an assertion about every matrix-Cauchy ensemble.

For a commuting Hermitian \(d\)-tuple \(A\), choose an independent isotropic \(d\)-dimensional Cauchy vector \(Z\) with density
\[
 p_\varepsilon(z)=\frac{\Gamma((d+1)/2)}{\pi^{(d+1)/2}}
 \frac{\varepsilon}{(\varepsilon^2+|z|^2)^{(d+1)/2}}.
\]
The aid's joint-functional-calculus extension is
\[
 \mathbb E_Z f(A_1+Z_1I,\ldots,A_d+Z_dI)
 =(p_\varepsilon*f)(A_1,\ldots,A_d),
\]
conditional on \(A\), with derivative bounds
\[
 \|\partial^\gamma(p_\varepsilon*f)\|_\infty
 \le C_{d,\gamma}\varepsilon^{-|\gamma|}\|f\|_\infty.
\]
The derivative exponent is scale sharp. This joint Cauchy density is not generally a product of independent scalar Cauchy laws. For noncommuting tuples, an arbitrary scalar function on \(\R^d\) has no specified joint functional calculus; moreover Cauchy moments prevent an unrestricted finite-expectation polynomial extension. These are obstructions to a naive extension, not a theorem excluding every noncommutative smoothing construction.

For the second question, take self-adjoint \(x_1,\ldots,x_m\) in a tracial von Neumann algebra, with faithful normal trace, and work in \(\mathcal N=W^*(x_1,\ldots,x_m)\). Assume the free difference quotients on the polynomial algebra are well-defined densely defined operators
\[
 \partial_i x_j=\delta_{ij}1\otimes1,\qquad
 J_i=\partial_i^*(1\otimes1)\in L^2(\mathcal N,\tau).
\]
Existence of the adjoint values is a hypothesis. Put \(D_i=m_{\mathrm{flip}}\partial_i\), where \(m_{\mathrm{flip}}(a\otimes b)=ba\), and \(D=(D_1,\ldots,D_m)\). The source asks whether there are common polynomials \(H_k\) such that
\[
 \sum_i\|D_iH_k(x)-J_i\|_2^2\longrightarrow0,
 \quad\text{that is }J\in\overline{\operatorname{Ran}D}^{\,L^2(\mathcal N)^m}.
\]
It is not enough to approximate each coordinate using unrelated potentials.

The aid's exact Hilbert-space test is
\[
 J\in\overline{\operatorname{Ran}D}
 \quad\Longleftrightarrow\quad
 \langle J,F\rangle=0\text{ for every }F\in\ker D^*,
 \qquad
 \operatorname{dist}(J,\overline{\operatorname{Ran}D})
 =\|P_{\ker D^*}J\|_2.
\]
The kernel uses the actual adjoint domain; this criterion does not itself prove annihilation in the general problem. In one variable, polynomial antiderivatives make the cyclic-gradient range dense in \(L^2(\mathcal N)\), so the assertion holds whenever the conjugate variable exists. For a centered nondegenerate correlated semicircular family with covariance matrix \(\Sigma\succ0\), it holds exactly with
\[
 J_i=\sum_j(\Sigma^{-1})_{ij}x_j,\qquad
 H(X)=\tfrac12\sum_{i,j}(\Sigma^{-1})_{ij}X_iX_j,
 \qquad D_iH(x)=J_i.
\]
No general multivariable closure theorem is claimed from these special cases.
\subsection{Short English statement}
The fused record contains two questions. The aid gives exact multivariable Poisson smoothing for commuting tuples using central Cauchy noise, and an adjoint-kernel criterion for cyclic-gradient closure, with one-variable and correlated-semicircular positive cases. Neither is a general noncommutative solution.
\subsection{Sources}
\begin{itemize}
\item Local registry, permanent UnsolvedMath ID \(20002664\), exact public code \texttt{catalogue entry}; canonical location \texttt{aim-probability-notes.json}, \texttt{notes[105]}.
\item AIM, \emph{Free Analysis}, 2006 workshop questions; supplied reference \url{https://aimath.org/WWN/freeanalysis/freeanalysis.pdf}. The verification field says the first question ends on page 1 and the second starts on page 2, with page number 1 attached to ``Q:''; it identifies U+000F as corrupted epsilon and the arrow extraction as a corrupted maps-to symbol. No PDF was accessed here.
\item Supplied research attempt 1: \texttt{partial\_result}, run status \texttt{open}, rigor label \texttt{proved}, low novelty confidence; registry status \texttt{partially\_solved}, classification \texttt{PARTIAL-PROGRESS}, statement status \texttt{reconstructed\_unverified}, created and updated 14 August 2026. No proposer, original remarks, literature entry, or primary citations for the supplied partial results are given.
\item The aid is machine generated, not peer reviewed, and explicitly requires independent verification before citation.
\end{itemize}
\end{document}
